% The Torus as Group and as Surface  (source §3.3–3.5; added: full S^1 x S^1
% proof, genus / Euler-characteristic computation via a CW structure, and the
% distinction between the topological and the complex structure)
\chapter{The Torus as Group and as Surface}
\label{ch:torus_group_surface}

\Cref{ch:lattices} made \(T = \C/\Lambda\) into a compact Riemann surface. This
chapter adds the two structures that make it the right home for the elliptic
group law. First, \(T\) is an abelian group, and its operations are holomorphic,
so \(T\) is a complex Lie group; this is the group whose law \cref{thm:main}
transports to the cubic. Second, \(T\) is topologically a torus of genus one,
and that single integer is what gives Riemann--Roch its simple shape in
\cref{ch:riemann_roch}. We close by separating what is topological (the same for
every \(\Lambda\)) from what is complex-analytic (the part that depends on
\(\Lambda\)).

\section{The Group Law}
\label{sec:tg-group-law}

The complex plane is an abelian group under addition, and \(\Lambda\) is a
subgroup, so the quotient inherits an operation.

\begin{proposition}
\label{prop:tg-welldef}
The formula
\[
  [u] + [v] = [u + v]
\]
is well defined on \(\C/\Lambda\) and makes it an abelian group with identity
\(0_T = [0]\) and inverses \(-[z] = [-z]\).
\end{proposition}

\begin{proof}
Suppose \(u' \equiv u\) and \(v' \equiv v \pmod \Lambda\), say
\(u' = u + \lambda_1\) and \(v' = v + \lambda_2\) with
\(\lambda_1, \lambda_2 \in \Lambda\). Then
\(u' + v' = (u + v) + (\lambda_1 + \lambda_2)\), and
\(\lambda_1 + \lambda_2 \in \Lambda\) because \(\Lambda\) is a subgroup, so
\([u' + v'] = [u + v]\); the operation does not depend on representatives. The
group axioms are inherited from \((\C, +)\): associativity and commutativity
descend immediately, \([0]\) is a two-sided identity, and \([-z]\) is the
inverse of \([z]\).
\end{proof}

Equivalently, \(\pi : \C \to T\) is a surjective homomorphism of abelian groups
with kernel \(\Lambda\), so \(T \cong \C/\Lambda\) as groups --- the quotient in
the group-theoretic sense agrees with the quotient as a space. Adding two
classes is carried out by lifting, adding in \(\C\), and reducing back into the
fundamental parallelogram (\cref{fig:grouplaw}).

\begin{figure}[ht]
\centering
\begin{tikzpicture}[scale=1.05, every node/.style={font=\small}]
  % Fundamental parallelogram with sage-verified geometry
  \draw[draw=formteal!55!black, fill=formteal!8, line width=0.7pt]
    (0,0) -- (\latwax,\latway) -- (\latsxc,\latsyc) -- (\latwbx,\latwby) -- cycle;
  % Adjacent translate (where u+v lands)
  \draw[draw=formteal!45!black, dashed, line width=0.35pt, opacity=0.5]
    (\latwax,\latway) -- (2*\latwax,2*\latway)
    -- ({2*\latwax+\latwbx},{2*\latway+\latwby})
    -- (\latsxc,\latsyc);
  % Two points u, v
  \fill[formrust!85!black] (\grupux,\grupuy) circle (1.8pt);
  \node[formrust!60!black, below right] at (\grupux,\grupuy) {$[u]$};
  \fill[formrust!85!black] (\grupvx,\grupvy) circle (1.8pt);
  \node[formrust!60!black, above left] at (\grupvx,\grupvy) {$[v]$};
  % u+v lifted
  \fill[formgold!85!black] (\grupsumx,\grupsumy) circle (1.8pt);
  \node[formgold!50!black, above right] at (\grupsumx,\grupsumy) {$u+v$};
  \draw[-{Latex[length=2mm]}, draw=formgold!75!black, line width=0.7pt, dashed]
    (\grupux,\grupuy) -- (\grupsumx,\grupsumy);
  \draw[-{Latex[length=2mm]}, draw=formgold!75!black, line width=0.7pt, dashed]
    (\grupvx,\grupvy) -- (\grupsumx,\grupsumy);
  % Reduced representative (inside P_Lambda) via subtracting omega_1
  \fill[formplum] (\grupredx,\grupredy) circle (1.8pt);
  \node[formplum, above left, font=\small] at (\grupredx,\grupredy) {$[u+v]$};
  \draw[-{Latex[length=2.5mm]}, draw=formplum, line width=0.9pt]
    (\grupsumx,\grupsumy) to[bend left=22] (\grupredx,\grupredy);
  \node[formplum, font=\scriptsize] at (1.7,2.6) {translate by $-\omega_1$};
  % Origin
  \fill[black] (0,0) circle (1.6pt) node[below left] {$[0]$};
  % Annotation with sage-verified arithmetic
  \node[draw=black!40, rounded corners=2pt, fill=formivory, align=left, anchor=west,
        inner sep=4pt, font=\scriptsize] at (4.5,1.45)
    {$u=0.40\,\omega_1+0.30\,\omega_2$\\
     $v=0.70\,\omega_1+0.60\,\omega_2$\\
     $u+v=1.10\,\omega_1+0.90\,\omega_2$\\
     reduce: subtract $\omega_1$\\
     to get $[u+v]=0.10\,\omega_1+0.90\,\omega_2$};
\end{tikzpicture}
\caption{Addition on \(T = \C/\Lambda\): lift \([u]\) and \([v]\) to \(\C\), add
to get \(u+v\), then reduce modulo \(\Lambda\) back into \(P_\Lambda\). The
arithmetic shown is computed by \texttt{sagemath/02\_torus\_group\_law.sage}.}
\label{fig:grouplaw}
\end{figure}

The group law is not merely set-theoretic: it is compatible with the complex
structure of \cref{ch:lattices}.

\begin{proposition}[\texorpdfstring{$T$}{T} is a complex Lie group]
\label{prop:tg-lie}
Addition \(m : T \times T \to T\), \(m([z],[w]) = [z+w]\), and negation
\(\iota : T \to T\), \(\iota([z]) = [-z]\), are holomorphic. Consequently every
translation
\[
  \tau_a : T \to T, \qquad \tau_a([z]) = [z + a] \quad (a \in \C)
\]
is a biholomorphic automorphism of \(T\), and the translations act simply
transitively on \(T\).
\end{proposition}

\begin{proof}
In the charts of \cref{thm:lat-riemann-surface} the maps \(m\) and \(\iota\)
lift to \((z,w) \mapsto z+w\) and \(z \mapsto -z\) on open subsets of \(\C\),
which are holomorphic; since \(\pi\) is a local biholomorphism, \(m\) and
\(\iota\) are holomorphic on \(T\). Fixing \(a\), the translation \(\tau_a\) is
holomorphic with holomorphic inverse \(\tau_{-a}\), hence biholomorphic. Given
\([z], [w] \in T\), the unique translation carrying \([z]\) to \([w]\) is
\(\tau_{w-z}\), so the action is simply transitive.
\end{proof}

\begin{remark}
\label{rem:tg-homogeneous}
Simple transitivity says \(T\) looks the same at every point: the translation
\(\tau_a\) is a biholomorphism carrying \(0_T\) to \([a]\). We use this
homogeneity constantly --- for instance to move a divisor so that one of its
points sits at the origin --- and it is the reason a single local computation at
\(0_T\) often suffices.
\end{remark}

\section{The Torus as a Product of Circles}
\label{sec:tg-product}

The abstract law \([u]+[v]=[u+v]\) has a concrete realization that makes
commutativity transparent.

\paragraph{The circle group.}
The unit circle
\[
  S^1 = \{ z \in \C : |z| = 1 \} = \{ e^{i\theta} : \theta \in \R \}
\]
is an abelian group under complex multiplication, since
\(e^{i\alpha} e^{i\beta} = e^{i(\alpha+\beta)} = e^{i\beta} e^{i\alpha}\). The
homomorphism \(t \mapsto e^{2\pi i t}\) from \(\R\) onto \(S^1\) has kernel
\(\Z\), so the first isomorphism theorem gives
\[
  \R/\Z \;\cong\; S^1 .
\]
The circle is the real line wound onto itself, with the group law becoming angle
addition.

\begin{proposition}[Torus \(\cong\) product of circles]
\label{prop:tg-s1xs1}
Let \(\Lambda = \Z\omega_1 + \Z\omega_2\). The map
\[
  \Psi : \C/\Lambda \longrightarrow S^1 \times S^1,
  \qquad
  [\,s\omega_1 + t\omega_2\,] \longmapsto \bigl( e^{2\pi i s},\, e^{2\pi i t} \bigr)
\]
is an isomorphism of topological groups.
\end{proposition}

\begin{proof}
Since \(\omega_1, \omega_2\) are \(\R\)-independent, every \(z \in \C\) has a
unique expression \(z = s\omega_1 + t\omega_2\) with \(s, t \in \R\). The
real-linear isomorphism \(\phi : \R^2 \to \C\), \((s,t) \mapsto s\omega_1 +
t\omega_2\), carries \(\Z^2\) onto \(\Lambda\), so it descends to an isomorphism
\(\R^2/\Z^2 \xrightarrow{\sim} \C/\Lambda\) of abelian groups. Composing the
inverse with the coordinatewise exponential
\(\R^2/\Z^2 = (\R/\Z)^2 \to S^1 \times S^1\) yields \(\Psi\).

\emph{Homomorphism.} Writing \([z_k] = [s_k\omega_1 + t_k\omega_2]\),
\[
  \Psi([z_1]+[z_2])
  = \bigl( e^{2\pi i(s_1+s_2)}, e^{2\pi i(t_1+t_2)} \bigr)
  = \bigl( e^{2\pi i s_1} e^{2\pi i s_2}, e^{2\pi i t_1} e^{2\pi i t_2} \bigr)
  = \Psi([z_1])\,\Psi([z_2]).
\]
\emph{Injectivity.} If \(\Psi([z]) = (1,1)\) then \(e^{2\pi i s} = e^{2\pi i t}
= 1\), so \(s, t \in \Z\) and \(z \in \Lambda\), i.e.\ \([z] = [0]\).
\emph{Surjectivity.} The class \([s\omega_1 + t\omega_2]\) maps to
\((e^{2\pi i s}, e^{2\pi i t})\). \emph{Homeomorphism.} \(\Psi\) is a continuous
bijection between compact Hausdorff spaces, hence a homeomorphism.
\end{proof}

\begin{corollary}
\label{cor:tg-abelian}
\(T = \C/\Lambda\) is an abelian group, and as a topological group it is
independent of \(\Lambda\): every complex torus is isomorphic to
\(S^1 \times S^1\).
\end{corollary}

\begin{proof}
\(S^1\) is abelian, a product of abelian groups is abelian, and \(\Psi\)
transports this to \(T\). The isomorphism class of \(S^1 \times S^1\) does not
involve \(\Lambda\).
\end{proof}

\begin{keyidea}{Why commutativity is obvious from the circle picture}
A class \([s\omega_1 + t\omega_2]\) corresponds to a pair of angles
\((e^{2\pi i s}, e^{2\pi i t})\), one winding in the \(\omega_1\)-direction and
one in the \(\omega_2\)-direction. Adding classes rotates the two circles
independently,
\[
  (e^{2\pi i s_1}, e^{2\pi i t_1})\,(e^{2\pi i s_2}, e^{2\pi i t_2})
  = (e^{2\pi i(s_1+s_2)}, e^{2\pi i(t_1+t_2)}),
\]
and rotating by \(\alpha\) then \(\beta\) equals rotating by \(\beta\) then
\(\alpha\). Commutativity of the group law is just commutativity of adding real
angles.
\end{keyidea}

\begin{warning}
\label{warn:tg-not-biholo}
\Cref{prop:tg-s1xs1} is an isomorphism of \emph{topological} groups, not of
Riemann surfaces. The product \(S^1 \times S^1\) carries no preferred complex
structure; the map \(\Psi\) uses the real coordinates \(s, t\) and is only
real-analytic, not holomorphic. The complex structure of \(T\), which remembers
the shape \(\tau\) of \cref{ch:lattices}, is exactly the information that
\(\Psi\) discards. We return to this in \cref{sec:tg-top-vs-complex}.
\end{warning}

\section{Genus and Euler Characteristic}
\label{sec:tg-genus}

Topologically \(T\) is built from the parallelogram \(P_\Lambda\) by gluing
opposite edges, and this gluing gives it a cell decomposition from which the
Euler characteristic is immediate.

\begin{proposition}
\label{prop:tg-euler}
The torus \(T\) has Euler characteristic \(\chi(T) = 0\), and hence genus
\(g = 1\).
\end{proposition}

\begin{proof}
Identifying opposite edges of \(P_\Lambda\) as in \cref{fig:cwtorus} produces a
CW structure on \(T\): the four corners are all glued to a single point, the two
pairs of opposite edges become two edge-loops \(a, b\), and the interior is one
\(2\)-cell. Thus
\[
  V = 1, \qquad E = 2, \qquad F = 1,
\]
so
\[
  \chi(T) = V - E + F = 1 - 2 + 1 = 0 .
\]
(As a check, \(\chi(S^1 \times S^1) = \chi(S^1)\,\chi(S^1) = 0 \cdot 0 = 0\) by
the product formula, agreeing with \cref{prop:tg-s1xs1}.) For a compact,
connected, orientable surface the Euler characteristic and the genus \(g\) --- the
number of handles --- are related by \(\chi = 2 - 2g\). Setting \(2 - 2g = 0\)
gives \(g = 1\).
\end{proof}

\begin{figure}[ht]
\centering
\begin{tikzpicture}[scale=1.2, every node/.style={font=\small}]
  % The identified square: bottom and top are edge a, left and right are edge b
  \coordinate (BL) at (0,0);
  \coordinate (BR) at (3,0);
  \coordinate (TR) at (3,3);
  \coordinate (TL) at (0,3);
  \fill[formteal!7] (BL) -- (BR) -- (TR) -- (TL) -- cycle;
  % edge a: bottom and top, single arrow
  \draw[postaction={decorate, decoration={markings, mark=at position 0.55 with {\arrow[draw=formrust!75!black]{>}}}}, formrust!75!black, line width=1pt] (BL) -- (BR);
  \draw[postaction={decorate, decoration={markings, mark=at position 0.55 with {\arrow[draw=formrust!75!black]{>}}}}, formrust!75!black, line width=1pt] (TL) -- (TR);
  % edge b: left and right, double arrow
  \draw[postaction={decorate, decoration={markings, mark=at position 0.5 with {\arrow[draw=formgold!75!black]{>}}, mark=at position 0.6 with {\arrow[draw=formgold!75!black]{>}}}}, formgold!75!black, line width=1pt] (BL) -- (TL);
  \draw[postaction={decorate, decoration={markings, mark=at position 0.5 with {\arrow[draw=formgold!75!black]{>}}, mark=at position 0.6 with {\arrow[draw=formgold!75!black]{>}}}}, formgold!75!black, line width=1pt] (BR) -- (TR);
  % all four corners are the single vertex v
  \foreach \p in {BL,BR,TR,TL}{\fill[formteal!80!black] (\p) circle (2pt);}
  \node[below left]  at (BL) {$v$};
  \node[below right] at (BR) {$v$};
  \node[above right] at (TR) {$v$};
  \node[above left]  at (TL) {$v$};
  \node[formrust!70!black] at (1.5,-0.3) {$a$};
  \node[formrust!70!black] at (1.5,3.3)  {$a$};
  \node[formgold!60!black] at (-0.3,1.5) {$b$};
  \node[formgold!60!black] at (3.3,1.5)  {$b$};
  \node[formteal!55!black] at (1.5,1.5) {$F$};
  % Count annotation
  \node[draw=black!40, rounded corners=2pt, fill=formivory, align=left,
        inner sep=4pt, font=\scriptsize, anchor=west] at (3.9,1.5)
    {$V = 1$ (all corners $= v$)\\
     $E = 2$ (loops $a, b$)\\
     $F = 1$ (the square)\\[2pt]
     $\chi = 1 - 2 + 1 = 0$\\
     $\chi = 2 - 2g \Rightarrow g = 1$};
\end{tikzpicture}
\caption{The CW structure on the torus. Gluing the two \(a\)-edges and the two
\(b\)-edges of the square sends all four corners to one vertex \(v\), leaves two
edge-loops \(a, b\), and one \(2\)-cell \(F\); hence \(\chi(T) = 0\) and
\(g = 1\).}
\label{fig:cwtorus}
\end{figure}

\begin{remark}
\label{rem:tg-genus-matters}
The genus is the one topological number that controls the function theory of a
compact Riemann surface. For genus one the Riemann--Roch theorem
(\cref{ch:riemann_roch}) takes the especially clean form \(\ell(D) = \deg D\) for
divisors of positive degree, and it is this that pins down the dimensions of the
spaces of functions used to build \(\wp\) and \(\wp'\). Everything downstream
rests on \(g = 1\).
\end{remark}

\section{Topological versus Complex Structure}
\label{sec:tg-top-vs-complex}

It is worth separating clearly the two layers we have placed on \(T\).

\begin{itemize}
\item \textbf{Topological / smooth layer (the same for all \(\Lambda\)).} By
\cref{cor:tg-abelian} every complex torus is, as a topological group, the single
object \(S^1 \times S^1\); by \cref{prop:tg-euler} it has genus one. No invariant
at this level distinguishes one lattice from another --- all complex tori are
diffeomorphic.

\item \textbf{Complex-analytic layer (depends on \(\Lambda\)).} The holomorphic
structure of \cref{thm:lat-riemann-surface} is extra data. By
\cref{prop:lat-normal-form} it is recorded, up to isomorphism, by the point
\(\tau \in \HH\) modulo \(\mathrm{SL}_2(\Z)\); two tori \(\C/\Lambda\) and
\(\C/\Lambda'\) are biholomorphic if and only if their \(\tau\)-classes agree.
\end{itemize}

Thus the diffeomorphism type is rigid (one torus) while the biholomorphism type
varies in the one-parameter family \(\HH / \mathrm{SL}_2(\Z)\). The single
complex number that separates the classes is the \(j\)-invariant of
\cref{ch:synthesis}. \Cref{warn:tg-not-biholo} is the concrete face of this: the
real coordinates \((s,t)\) trivialize \(T\) topologically but are blind to the
complex structure that the rest of the book exploits.

\section{Summary}
\label{sec:tg-summary}

The torus \(T = \C/\Lambda\) is an abelian group under \([u]+[v]=[u+v]\), with
holomorphic addition and negation, so it is a complex Lie group on which the
translations act simply transitively (\cref{prop:tg-welldef,prop:tg-lie}). As a
topological group it is \(S^1 \times S^1\), which makes commutativity visibly the
commutativity of adding angles, but this identification is only real-analytic and
forgets the complex structure (\cref{prop:tg-s1xs1,warn:tg-not-biholo}). A CW
decomposition gives \(\chi(T) = 0\) and genus \(g = 1\)
(\cref{prop:tg-euler}), the topological fact behind the simple form of
Riemann--Roch. All complex tori share this topology; what varies is the complex
structure, recorded by \(\tau \in \HH/\mathrm{SL}_2(\Z)\).

\section*{Exercises}
\addcontentsline{toc}{section}{Exercises}

\begin{exercise}
\label{ex:tg-lie-holo}
Complete the details of \cref{prop:tg-lie}: writing the addition map in two
overlapping charts of \cref{thm:lat-riemann-surface}, verify explicitly that the
transition makes \(m : T \times T \to T\) holomorphic, and deduce that
\(\iota([z]) = [-z]\) is biholomorphic.
\end{exercise}

\begin{exercise}
\label{ex:tg-two-torsion}
Show that the negation map \(\iota([z]) = [-z]\) has exactly four fixed points on
\(T\), namely
\[
  [0],\quad \Bigl[\tfrac{\omega_1}{2}\Bigr],\quad
  \Bigl[\tfrac{\omega_2}{2}\Bigr],\quad
  \Bigl[\tfrac{\omega_1+\omega_2}{2}\Bigr].
\]
These are the \emph{two-torsion} points, the solutions of \(2[z] = [0]\); they
reappear as the branch points of \(\wp\) in \cref{ch:embedding} and as the
two-torsion of the group law in \cref{ch:chord_tangent_cases}.
\end{exercise}

\begin{exercise}
\label{ex:tg-first-iso}
Fill in the first-isomorphism-theorem argument for \(\R/\Z \cong S^1\): check
that \(t \mapsto e^{2\pi i t}\) is a continuous surjective homomorphism with
kernel \(\Z\), and that the induced bijection \(\R/\Z \to S^1\) is a
homeomorphism (use compactness of \([0,1]\)).
\end{exercise}

\begin{exercise}
\label{ex:tg-euler-refine}
Recompute \(\chi(T)\) using a finer CW structure --- for instance subdivide
\(P_\Lambda\) into two triangles before gluing --- and confirm you again obtain
\(\chi(T) = 0\), illustrating the invariance of the Euler characteristic.
\end{exercise}

\begin{exercise}
\label{ex:tg-homeo-not-biholo}
Show that \(\C/(\Z + \Z i)\) and \(\C/(\Z + \Z\,2i)\) are homeomorphic (indeed
diffeomorphic) as real surfaces. Granting the standard fact that the points
\(\tau = i\) and \(\tau = 2i\) lie in the closed fundamental domain
\(\{\,|\tau| \ge 1,\ |\operatorname{Re}\tau| \le \tfrac12\,\}\) of
\(\mathrm{SL}_2(\Z)\) and are not identified by it, conclude (using
\cref{ch:lattices}) that the two tori are \emph{not} biholomorphic. Explain why
this does not contradict \cref{cor:tg-abelian}.
\end{exercise}
